\scanpage{63}
% archive sheet number: 62
et tout chemin \inserted{$c$} définit des iso. d'ens à deux éléments
\diagram{\begin{tikzcd}[ampersand replacement=\&,column sep=large,row sep=small]
\underline{\omega}(\deleted{\unsure{D}}F)\simeq\underline{\omega}(D) \arrow[r,"{\underline{\omega}_X(c)}"] \& \underline{\omega}(\deleted{\unsure{D}}F')\simeq\underline{\omega}(D')\\
\unsure{\varepsilon_s|\underline{\varepsilon}_s}(F)=\varepsilon_s(D) \arrow[r] \& \varepsilon_s(F')\simeq\varepsilon_s(D')\\
\varepsilon_{a}(F)\simeq\varepsilon_a(D)\arrow[r] \& \varepsilon_a(F')\simeq\varepsilon_a(D')
\end{tikzcd}}

Ces ens. ne sont pas indépendants, chaque paire de deux détermine le troisième, par la relation
\[
\underline{\omega}(F)\simeq\varepsilon_s(F)\wedge\varepsilon_a(F)
\]
et les deux relations similaires équivalentes. Au niveau des caractères correspondants sur $\pi_1(\mathscr{X},F)$, cela donne
\[
\begin{aligned}
&\boxed{\chi_X(F)=\chi_s(F)\chi_a(F)}\\
&\overset{\shortparallel}{\underset{\underline{\omega}(D)}{\underset{\shortparallel}{\chi_X(D)}}}=\chi_s(\overset{\shortparallel}{D})\chi_a(D)
\end{aligned}
\]
ou de façon plus symétrique
\[
\chi_X(F)\rho_s(F)\chi_a(F)=1
\]

En plus de ces trois caractères, il y en a deux autres, $\chi_{\mathscr X}$ = caractère d'\unsure{orientabilité|orientation} de $\mathscr X$ (ne dépendant pas de la structure cellulaire), et $\chi_c$ (``combinatoire''), qui est donné par la parité \inserted{des côtés}\inserted{Attention, dans le \unsure{compte des points|points|parité}, \deleted{\unsure{du poids|est fixé}} il faut compter pour \uline{deux} le franchissement de \unsure{tout|d'un|d'une} côté except.$^{\mathrm{el}}$ (qui \unsure{au schéma|au schème|amorce un demi-tour}, se décompose en effet en deux opérations)}.
% slanted
\deleted{Le système local associé est les} \inserted{Le correspondant est défini par les} $\underline{\omega}(F)$, mais en $y$ affectant les transitions \uline{opposées} à celles données par \deleted{\unsure{identifications}} le transport parallèle sur $X$ défini par les $\underline{\omega}_X(c)$ -- donc on aura
\[
\underline{\omega}_{\mathscr X}(c)=\mathrm{par}(c)\,\omega_X(c)
\]
où $\mathrm{par}(c)$ désigne la parité de $c$.

\deleted{\unsure{En résumé|Le système}} Le système local associé à $\chi_c$ est \deleted{formé de} comme fibres les ens standards $\{\pm1\}$, mais les morphismes de transition donnés par les $\mathrm{par}(c)$. On a donc
\[
\boxed{\chi_{\mathscr X}=\chi_o\chi_c}.
\]
% slanted
\margin{On a $\chi_c=1$ cf page suivante !}

En résumé, compte tenu de ces relations, on dispose de trois caractères fondamentaux \inserted{parmi les cinq}, je choisis
\[
\chi_X,\chi_s,\chi_c
\]
donc les deux autres sont
\[
\chi_a=\chi_X\chi_s\quad\text{et surtout}\quad\chi_{\mathscr X}=\chi_X\chi_c.
\]

Il y a lieu d'expliciter les valeurs des caractères \unsure{principaux|précéd.|précédents}, si faire se peut, ainsi que celles des la Représentation
\[
u:\pi_1(\mathscr X,F_D)\longrightarrow\operatorname{Aut}(\mathrm{Pol}(D)),
\]
